\section{The f\/irst Painlev\'e equation and isomonodromy system}\label{section2}


Let us consider the {\em first Painlev\'e equation}
with a formal parameter $\hbar$:
\begin{gather*}
\PI\colon \
\hbar^{2} \frac{d^{2}q}{dt^{2}} = 6q^{2} + t.
\end{gather*}
The equation $\PI$ is obtained from
\begin{gather*}
\frac{d^{2} \tilde{q}}{d\tilde{t}^{2}} =
6 \tilde{q}^{2} + \tilde{t}
\end{gather*}
via the rescaling
$\tilde{t} = \hbar^{-4/5} t$,
$\tilde{q} = \hbar^{-2/5} q$.
We will regard $\hbar$ as a small parameter (i.e., Planck's constant),
and investigate a particular formal solution
of $\PI$ which has an $\hbar$-expansion.

\subsection[Formal solution of $\PI$]{Formal solution of $\boldsymbol{\PI}$}

$\PI$ has the following formal power series solution:
\begin{gather} \label{eq:formal-q}
q(t,\hbar) = \sum_{n=0}^{\infty} \hbar^{2n}q_{2n}(t)
= q_{0}(t) + \hbar^{2} q_{2}(t)
+ \hbar^{4} q_{4}(t) + \cdots.
\end{gather}
It contains only even order terms of $\hbar$ since $\PI$
is invariant under $\hbar \mapsto - \hbar$.
The leading term $q_{0} = q_{0}(t)$ satisf\/ies
\begin{gather} \label{eq:q0}
6q_{0}^{2} + t = 0, \qquad \text{hence}~q_{0}(t)=\sqrt{-t/6},
\end{gather}
and the subleading terms are recursively determined by
\begin{gather} \label{eq:recursion-q-2k}
q_{2(k+1)}(t) = \frac{1}{12 q_{0}(t)} \left(
\frac{d^{2} q_{2k}}{dt^{2}}(t)  -
6 \sum_{k_{1}+k_{2} = k+1,\,  k_{i}>0} q_{2k_{1}}(t) q_{2k_{2}}(t)
\right), \qquad k \ge 1.
\end{gather}
As we will see, the coef\/f\/icients of the formal series appearing in
this paper are multivalued functions of $t$ and are def\/ined
on the Riemann surface of~$q_{0}$. Thus, in what follows,
we may use~$q_{0}$ instead of~$t$ when we express coef\/f\/icients.

The relation \eqref{eq:recursion-q-2k} implies
\begin{gather*}
q_{2k} =
c_{2k} q_{0}^{1-5k}, \qquad c_{2k} \in \bbC .
\end{gather*}
It is obvious that the coef\/f\/icients $q_{2k}(t)$ have
a singularity at $q_{0} = 0$ (i.e., $t = 0$).
This special point is called a {\em turning point} of $\PI$
\cite[Def\/inition 2.1]{KT-PI} (see also \cite[Section~4]{KT05}).
Throughout the paper, we assume the following:

\begin{ass}
The independent variable $t$ of $\PI$
lies on a domain that doesn't contain the origin.
\end{ass}


\begin{rem}
The formal solution \eqref{eq:formal-q} is called a
{\em $0$-parameter solution} of $\PI$ in \cite{KT05}
since it doesn't contain free parameters.
More general formal solutions having
one or two free para\-me\-ters (called 1- or 2-parameter solutions)
are constructed in~\cite{AKT-P} for all Painlev\'e equations
of second order.
See also~\cite{Umeta} for a construction of general formal
solutions of {\em higher} order Painlev\'e equations.
\end{rem}

\begin{rem}\sloppy
The formal solution \eqref{eq:formal-q} is in fact a~{\em divergent} series.
However, \cite[Theo\-rem~1.1]{Kamimoto-Koike} proved that
the formal solution is {\em Borel summable} when
$q_{0}$ satisf\/ies $q_{0} \ne 0$ and
\mbox{$\arg q_{0} \notin \{ \frac{2\ell}{5}\pi \,|\,
\ell \in {\mathbb Z} \}$}.
The exceptional set is called the {\em Stokes curve} of $\PI$.
(See \cite[Def\/ini\-tion~2.1]{KT-PI} for the notion of
Stokes curves of Painlev\'e equations
with a small parameter~$\hbar$.)
That is, there exists a function
which is analytic in $\hbar$ on a sectorial domain
with the center at the origin (which is also analytic in $t$)
such that~\eqref{eq:formal-q} is the asymptotic expansion
of the function for $\hbar \rightarrow 0$ in the sector.
The analytic function is called the {\em Borel sum}
of the formal series~\eqref{eq:formal-q}, and
it gives an analytic solution of $\PI$
(see~\cite{Costin08} for Borel summation method).
This particular asymptotic solution obtained by
the Borel summation method is called
the {\em tri-tronqu\'ee solution}
of $\PI$ (see~\cite{Joshi-Kitaev}), and the
{\em non-linear Stokes phenomena} on Stokes curves
are analyzed by \cite{FIKN, Kap, Takei}.
\end{rem}

\subsection[Isomonodromy system and the $\tau$-function]{Isomonodromy system and the $\boldsymbol{\tau}$-function}

It is known that $\PI$ describes the compatibility condition
for the following system of linear PDEs (cf.~\cite[Appendix C]{JM2}):
\begin{gather} \label{eq:Lax-matrix}
\hbar \frac{\p \Psi}{\p x}  = A \Psi, \qquad
\hbar \frac{\p \Psi}{\p t}  = B \Psi,
\end{gather}
where
\begin{gather*}
A   =
\begin{pmatrix}
A_{11} & A_{12} \\
A_{21} & A_{22}
\end{pmatrix}
:=
\begin{pmatrix}
  p
&
  4(x-q)
\\
\displaystyle x^{2} + q x + q^{2} + \frac{t}{2}
&
  - p
\end{pmatrix},
\\
B  =
\begin{pmatrix}
B_{11} & B_{12} \\
B_{21} & B_{22}
\end{pmatrix}
:=
\begin{pmatrix}
 0
&
   2
\\
\displaystyle \frac{x}{2} + q
&
  0
\end{pmatrix}.
\end{gather*}
The compatibility condition
\begin{gather*}
\hbar \frac{\p A}{\p t} - \hbar \frac{\p B}{\p x} + [A,B] = 0
\end{gather*}
is equivalent to the following Hamiltonian system
\begin{gather} \label{eq:Ham-PI}
\hbar \frac{dq}{dt} = \frac{\p H}{\p p}, \qquad
\hbar \frac{dp}{dt} = - \frac{\p H}{\p q},
\end{gather}
where the (time-dependent) Hamiltonian is given by
\begin{gather*}% \label{eq:Ham}
H = H(q,p,t) := \frac{1}{2}p^{2} - 2 q^{3} - t q.
\end{gather*}
We can easily check that \eqref{eq:Ham-PI} and $\PI$ are equivalent.
The above system of linear ODEs
is called the {\em isomonodromy system} associated with $\PI$
(see~\cite{JMU, JM2}).

Let $(q,p) = (q(t,\hbar), p(t,\hbar))$ be a formal power series
solution of the Hamiltonian system~\eqref{eq:Ham-PI}; that is,
$q(t,\hbar)$ is the formal solution~\eqref{eq:formal-q}
of~$\PI$, and
\begin{gather*} %\label{eq:formal-p}
p(t,\hbar) = \hbar \frac{d q(t,\hbar)}{dt} =
\sum_{n=0}^{\infty} \hbar^{2n+1} p_{2n+1}(t).
\end{gather*}
The corresponding Hamiltonian function is denoted by
\begin{gather} \label{eq:sigma}
\sigma(t,\hbar) := H\big(q(t,\hbar), p(t,\hbar), t\big).
\end{gather}
We can check that \eqref{eq:sigma}
is invariant under $\hbar \mapsto - \hbar$, and hence
it has the following expansion:
\begin{gather} \label{eq:parity-sigma}
\sigma(t,\hbar) = \sum_{n=0}^{\infty} \hbar^{2n} \sigma_{2n}(t).
\end{gather}

\begin{Def} [\cite{JMU, Okamoto}] \label{def:tau-function}
The {\em $\tau$-function}
(corresponding to the formal solution \eqref{eq:formal-q})
of $\PI$ is def\/ined by
\begin{gather} \label{eq:tau}
\hbar^{2} \frac{d}{dt} \log \tau(t,\hbar) = \sigma(t,\hbar)
\end{gather}
up to constant.
\end{Def}

The $\tau$-function can also be def\/ined in terms of
a solution of~\eqref{eq:Lax-matrix}~\cite{JMU} (see also
Appendix~\ref{section:relation-of-two-conjectures}).
The expansion~\eqref{eq:parity-sigma} implies that the $\tau$-function~\eqref{eq:tau} has an expansion of the form
\begin{gather*} % \begin{equation}
\log \tau(t,\hbar) = \sum_{g=0}^{\infty} \hbar^{2g-2} \tau_{2g}(t).
\end{gather*} % \end{equation}

\subsection{Spectral curve}

In what follows, we assume that the formal solution
$(q(t,\hbar), p(t,\hbar))$ of~\eqref{eq:Ham-PI}
constructed above is substituted into the coef\/f\/icients
of the isomonodromy system \eqref{eq:Lax-matrix}.
Then, the coef\/f\/icients of the isomonodromy system has the following
$\hbar$-expansions:
\begin{gather*}
A   =   A_{0}(x,t) + \hbar A_{1}(x,t) + \hbar^{2} A_{2}(x,t) + \cdots, \\
B   =   B_{0}(x,t) + \hbar B_{1}(x,t) + \hbar^{2} B_{2}(x,t) + \cdots,
\end{gather*}
whose top terms are given by
\begin{gather*}
A_{0}(x,t)   =
\begin{pmatrix}
  0
&
 4(x-q_{0})
\\
\displaystyle x^{2} + q_{0} x + q_{0}^{2} + \frac{t}{2}
&
  0
\end{pmatrix}, \qquad
%\label{eq:B0-matrix}
B_{0}(x,t)    =
\begin{pmatrix}
  0
&
  2
\\
\displaystyle \frac{x}{2}+q_{0}
&
   0
\end{pmatrix}.
\end{gather*}
Observe that, since $q_{0}$ satisf\/ies \eqref{eq:q0},
the algebraic curve def\/ined by
\begin{gather} \label{eq:sp-curve}
\det  (y - A_{0}(x,t)  ) =
y^{2} - 4  ( x-q_{0}  )^{2}  ( x+2q_{0}  )
= 0
\end{gather}
has {\em genus $0$}. Actually, this gives a family of algebraic
curves in ${\mathbb C}^{2}_{(x,y)}$
parametrized by $t$. Since we have assumed that
$t \ne 0$, $x=q_{0}$ and $x=-2q_{0}$ are distinct.

\begin{Def}
We call the algebraic curve~\eqref{eq:sp-curve}
{\em the semi-classical spectral curve}, or
{\em the spectral curve}
of (the f\/irst equation of) the isomonodromy system~\eqref{eq:Lax-matrix}.
\end{Def}

\begin{rem}
It is shown in \cite[Proposition~1.3]{KT-PI} that, for all (second order)
Painlev\'e equations with a~formal parameter $\hbar$,
the semi-classical spectral curves corresponding to
the same type of formal power series solution
as \eqref{eq:formal-q} have {\em genus $0$}.
\end{rem}

\begin{rem}
Since we are taking the semi-classical limit
(i.e., top term in $\hbar$-expansion),
our spectral curve \eqref{eq:sp-curve} is dif\/ferent from
usual spectral curves for isomonodromic deformation equations
discussed, e.g., in~\cite{Olsha, Takasaki}.
The spectral curves in the above papers have {\em higher genus}.
Recently, Nakamura~\cite{Nakamura} investigates
the geometry of genus~$2$ spectral curves which appear
in an {\em autonomous limit} of the
{\em 4th order Painlev\'e equations},
and use them to classify the Painlev\'e equations.
See \cite{KNS} for the list of 4th order Painlev\'e equations.
\end{rem}

\subsection{WKB analysis of isomonodromy system in scalar form}
\label{subsection:WKB-for-scalar-Lax}

Denote the unknown vector function of \eqref{eq:Lax-matrix}
by $\Psi = {}^t(\psi_1, \psi_2)$.
Then, $\psi = \psi_1$ satisf\/ies the
following scalar version of isomonodromy system
\begin{gather}
\left( \left(\hbar\frac{\p}{\p x}\right)^{2} +
f \left(\hbar\frac{\p}{\p x}\right) + g \right)\psi = 0, \nonumber\\
\hbar \frac{\p \psi}{\p t} = \frac{1}{2(x-q)} \left(
\hbar \frac{\p \psi}{\p x} - p \psi \right),\label{eq:Lax-scalar}
\end{gather}
where
\begin{gather*}
f = f(x,t,\hbar)   :=   -\operatorname{tr} A - \hbar \frac{\p}{\p x}\log A_{12}
= -\hbar \frac{1}{x-q},  \\
g = g(x,t,\hbar)   :=   \det A - \hbar \frac{\p A_{11}}{\p x} +
\hbar A_{11} \frac{\p}{\p x}\log A_{12}   \\
\hphantom{g}{}
  =
- \big(4x^{3} + 2tx + p^{2} - 4q^{3} - 2tq\big) + \hbar \frac{p}{x-q}.
\end{gather*}
The coef\/f\/icients of $f$ and $g$ have an $\hbar$-expansion since
$q$ and $p$ are contained in them
\begin{gather} \label{eq:hbar-correction-in-quantization}
f   =   - \hbar  \frac{1}{x-q_{0}}
 +\hbar^{3} \frac{1}{1728q_{0}^{4}(x-q_{0})^{2}} + \hbar^{5}
 \frac{49x-51q_{0}}{5971968q_{0}^{9}(x-q_{0})^{3}} + \cdots,
\\  \label{eq:hbar-correction-in-quantization2}
g   =   -4(x-q_{0})^{2}(x+2q_{0})
 - \hbar^{2} \frac{x+11q_{0}}{144q_{0}^{2}(x-q_{0})}
 - \hbar^{4} \frac{7x^{2}+34q_{0}x-53q_{0}^{2}}
  {248832q_{0}^{7}(x-q_{0})^{2}}+\cdots .
\end{gather}
The top term of $g$ appears in the def\/ining equation of
the spectral cur\-ve~\eqref{eq:sp-curve},
and its zeros are called {\em turning points}
of the f\/irst equation of~\eqref{eq:Lax-scalar}
in the WKB analysis.
In particular, under the assumption $t \ne 0$, there is
\begin{itemize}\itemsep=0pt
\item
a {\em simple} turning point at $x=-2q_{0}$
which is a branch point
of the spectral curve \eqref{eq:sp-curve}, and
\item
a {\em double} turning point at $x = q_{0}$
which is a singular point of
the spectral curve \eqref{eq:sp-curve}.
\end{itemize}
Consider the {\em Riccati equation}
\begin{gather} \label{eq:Riccati}
\hbar^{2} \left( P^{2} + \frac{\p P}{\p x} \right)
+ f \hbar P + g = 0.
\end{gather}
This is equivalent to the f\/irst equation
in \eqref{eq:Lax-scalar} by
\begin{gather*}
\psi = \exp\left( \int^{x} Pdx \right), \qquad
 \text{i.e.}, \quad P = \frac{1}{\psi} \frac{\p \psi}{\p x}  .
\end{gather*}
Let
\begin{gather*}
P^{(\pm)}(x,t,\hbar) = \sum_{m=0}^{\infty}\hbar^{m-1}P^{(\pm)}_{m}(x,t)
\end{gather*}
be the formal solutions of \eqref{eq:Riccati} with the top term
\begin{gather*} %\label{eq:P0-pm}
P^{(\pm)}_{0}(x,t) = \pm 2(x-q_{0})\sqrt{x+2q_{0}}.
\end{gather*}
The coef\/f\/icients $P^{(\pm)}_{m}(x,t)$ are recursively determined by
\begin{gather} \label{eq:recursion-for-Pm}
2P^{(\pm)}_{0}P^{(\pm)}_{m+1} +
\sum_{\substack{a+b=m+1 \\ a,b \ge 1}}
P^{(\pm)}_{a} \big(P^{(\pm)}_{b} + f_{b}\big) +
\frac{\p P^{(\pm)}_{m}}{\p x} + g_{m+1} = 0 \qquad \text{for} \quad m \ge 0,
\end{gather}
where $f_{a}$ and $g_{a}$ are the coef\/f\/icient of $\hbar^{a}$
in $f$ and $g$, respectively. Explicit forms of the f\/irst few terms
are given by
\begin{gather*}
P^{(\pm)}_{1}   =   -\frac{1}{4(x+2q_{0})}, \qquad
P^{(\pm)}_{2}   =   \pm \frac{x+17q_{0}}{576q_{0}^{2}(x+2q_{0})^{5/2}}, \qquad
P^{(\pm)}_{3}   =   -\frac{2x^{2}+20q_{0}x+77q_{0}^{2}}
{6912q_{0}^{4}(x+2q_{0})^{4}}, \\
P^{(\pm)}_{4}   =   \pm \frac{28x^{4}+500q_{0}x^{3}+
3684q_{0}^{2}x^{2}+14273q_{0}^{3}x+27307q_{0}^{4}}
{3981312 q_{0}^{7}(x+2q_{0})^{11/2}}.
\end{gather*}
It is obvious from~\eqref{eq:recursion-for-Pm}
that $P^{(\pm)}_{m}(x,t)$ are holomorphic except
at the turning points and $x = \infty$
(and multivalued for even $m$).
It also follows from the recursion relation~\eqref{eq:recursion-for-Pm} that
\begin{gather} \label{eq:asymp-P}
P^{(\pm)}(x,t,\hbar)
= \pm \left(\frac{2}{\hbar}x^{3/2} + \frac{t}{2\hbar}x^{-1/2}
\mp \frac{1}{4}x^{-1} + \frac{\sigma(t,\hbar)}{2\hbar} x^{-3/2}
+ O\big(x^{-2}\big) \right)
\end{gather}
holds when $x \rightarrow \infty$.

\begin{rem} \label{rem:asymptotics}
We can check that $P^{(\pm)}_{m}(x,t)$'s
have the following asymptotic expansion for \mbox{$x \rightarrow \infty$}
\begin{gather*}
P^{(\pm)}_{0}(x,t)   =   \pm \left( 2x^{3/2} + \frac{t}{2}x^{-1/2}
+ O\big(x^{-3/2}\big) \right), \qquad
P^{(\pm)}_{1}(x,t)  =   -\frac{1}{4}x^{-1} + O\big(x^{-3/2}\big), \\
P^{(\pm)}_{m}(x,t)   =    O\big(x^{-3/2}\big) \qquad \text{for} \quad m \ge 2,
\end{gather*}
and we have \eqref{eq:asymp-P} after summing up
$\hbar^{m-1}P^{(\pm)}_{m}(x,t)$. Once you know that $P^{(\pm)}(x,t,\hbar)$
has an asymptotic expansion in this sense, subleading terms in
\eqref{eq:asymp-P} can be computed from
the Riccati equation \eqref{eq:Riccati}.
\end{rem}

Def\/ine
\begin{gather*}
P_{\rm odd}(x,t,\hbar)   :=   \frac{1}{2}
\big( P^{(+)}(x,t,\hbar) - P^{(-)}(x,t,\hbar) \big), \\
P_{\rm even}(x,t,\hbar)   :=   \frac{1}{2}
\big( P^{(+)}(x,t,\hbar) + P^{(-)}(x,t,\hbar) \big).
\end{gather*}
It is easy to check that (cf.~\cite[Section~2]{KT05})
\begin{gather} \label{eq:odd-even-relation}
P_{\rm even}(x,t,\hbar) =
-\frac{1}{2}\frac{\p}{\p x}
\log \frac{ \hbar P_{\rm odd}(x,t,\hbar)}
{2(x-q(t,\hbar))}
\end{gather}
and
\begin{gather} \label{eq:parity}
P_{\rm odd}(\sigma(x),t,\hbar) = - P_{\rm odd}(x,t,\hbar)
\end{gather}
hold. Here $x$ is regarded as a coordinate on the spectral curve,
and $\sigma$ is the covering involution for the spectral curve:
$P^{(+)}(\sigma(x),t) = - P^{(-)}(x,t)$.

Since
\begin{gather*}
\frac{ \hbar P_{\rm odd}(x,t,\hbar)}{2(x-q(t,\hbar))} =
\sqrt{x+2q_{0}} \big( 1 + O(\hbar) \big),
\end{gather*}
the right hand-side of~\eqref{eq:odd-even-relation}
is the derivative of the formal power series
\begin{gather*}
-\frac{1}{2}\log\frac{ \hbar P_{\rm odd}(x,t,\hbar)}{2(x-q(t,\hbar))} =
-\frac{1}{4}\log(x+2q_{0}) + O(\hbar).
\end{gather*}
Thus the ambiguity of the branch of the logarithm
only appears in the top term, but
we care about the ambiguity since
it doesn't matter in our computation.


The following theorem was applied in the
{\it transformation theory of Painlev\'e equations}
in~\cite{KT-PI}. We will use the fact
in the proof of our main theorem.

\begin{thm}[\protect{cf.~\cite[Proposition~1.2 and Theorem~1.1]{KT-PI}}]
\label{thm:Podd-and-t-derivative} \quad
\begin{itemize}\itemsep=0pt
\item[$(i)$]
The formal series $P^{(\pm)}(x,t,\hbar)$ satisfies
\begin{gather} \label{eq:dt-P}
\hbar \frac{\p}{\p t}P^{(\pm)}(x,t,\hbar) = \frac{\p}{\p x}\left(
\frac{\hbar P^{(\pm)}(x,t,\hbar) - p(t,\hbar)}{2(x-q(t,\hbar))}
\right).
\end{gather}
In particular, $P_{\rm odd}(x,t,\hbar)$ satisfies
\begin{gather} \label{eq:dt-Podd}
\frac{\p}{\p t}P_{\rm odd}(x,t,\hbar) =
\frac{\p}{\p x}\left( \frac{P_{\rm odd}(x,t,\hbar)}{2(x-q(t,\hbar))}
\right).
\end{gather}

\item[$(ii)$]
All coefficients of $P^{(\pm)}(x,t,\hbar)$ are holomorphic
except at the simple turning point $x=-2q_{0}$ and $x= \infty$.
In particular, they are holomorphic at
the double turning point $x=q_{0}$.

\item[$(iii)$]
The formal series
\begin{gather}
\psi_{\pm}(x,t,\hbar)   :=    \exp\left(
\pm \int^{x}_{v} P_{\rm odd}(x',t,\hbar)dx'
- \frac{1}{2} \log \frac{ \hbar P_{\rm odd}(x,t,\hbar)}
{2(x-q(t,\hbar))}\right)  \nonumber \\
 \hphantom{\psi_{\pm}(x,t,\hbar)}{} =
\left(\frac{2(x-q(t,\hbar))}{\hbar P_{\rm odd}(x,t,\hbar)} \right)^{1/2}
\exp\left( \pm \int^{x}_{v} P_{\rm odd}(x',t,\hbar) dx' \right)\label{eq:WKB-IM}
\end{gather}
satisfies the isomonodoromy system~\eqref{eq:Lax-scalar}.
Here $v$ is the simple turning point $-2q_{0}$.
The integral from $v$ is defined by
\begin{gather} \label{eq:contour-integral-Podd}
\int^{x}_{v} P_{\rm odd}(x',t,\hbar) dx' =
\frac{1}{2} \int_{\gamma_x} P_{\rm odd}(x',t,\hbar) dx',
\end{gather}
where the path $\gamma_x$ is depicted in Fig.~{\rm \ref{fig:contour}}
$($cf.~{\rm \cite[Section~2]{KT05})}.
\end{itemize}
\end{thm}

\begin{figure}[t]\centering
 \includegraphics{Iwaki-Fig1}
 \caption{For a~given~$x$, the path~$\gamma_{x}$ starts from
 the point~$\sigma(x)$ and ends at~$x$.
 The wiggly lines designate a branch cut, and the solid (resp.\ dotted) part
 represents a~part of path on the f\/irst (resp.\ the second)
 sheet of the spectral curve.}
 \label{fig:contour}
 \end{figure}



\begin{proof}
Although the scalar version of isomonodromy system \eqref{eq:Lax-scalar}
is dif\/ferent from that used in~\cite{KT-PI}, they are
related by a gauge transformation $\psi \mapsto (x-q)^{1/2} \psi$.
Therefore, the equalities~\eqref{eq:dt-P} and~\eqref{eq:dt-Podd}
in~(i) together with the holomorphicity of each coef\/f\/icient
of $P_{\rm odd}(x,t,\hbar)$ at $x=q_{0}$ follows from
\cite[Proposition~1.2 and Theorem~1.1]{KT-PI}.
Then, it turns out that the coef\/f\/icients of
$P_{\rm odd}(x,t,\hbar)/(x-q(t,\hbar))$
are also holomorphic due to~\eqref{eq:dt-Podd}.
Then, \eqref{eq:odd-even-relation}
implies that each coef\/f\/icient of $P_{\rm even}(x,t,\hbar)$ is also
holomorphic at $x=q_{0}$.
Thus we have proved~(ii).

The claim (iii) follows from a straightforward computations
\begin{gather*}
\frac{1}{\psi_{\pm}} \frac{\p \psi_{\pm}}{\p x}  =
\pm P_{\rm odd} -\frac{1}{2}\frac{\p}{\p x} \log
\left(\frac{ \hbar P_{\rm odd}}{2(x-q)}\right)
= \pm P_{\rm odd} + P_{\rm even}
= P^{(\pm)}, \\
\hbar\frac{1}{\psi_{\pm}} \frac{\p \psi_{\pm}}{\p t}   =
\frac{1}{2}\left( \frac{-\hbar (dq/dt)}{x-q} - \frac{\hbar}{P_{\rm odd}}
\frac{\p P_{\rm odd}}{\p t} \right) \pm \int^{x}_{v}
\hbar\frac{\p P_{\rm odd}}{\p t} dx \\
\hphantom{\hbar\frac{1}{\psi_{\pm}} \frac{\p \psi_{\pm}}{\p t}  }{}
  =
-\frac{p}{2(x-q)} - \frac{\hbar}{P_{\rm odd}}
\left( \frac{1}{2(x-q)}\frac{\p P_{\rm odd}}{\p x} -
\frac{P_{\rm odd}}{2(x-q)^{2}} \right) \pm \hbar \frac{P_{\rm odd}}{2(x-q)} \\
\hphantom{\hbar\frac{1}{\psi_{\pm}} \frac{\p \psi_{\pm}}{\p t}  }{} =
-\frac{p}{2(x-q)} + \frac{\hbar P^{(\pm)}}{2(x-q)} =
\frac{1}{2(x-q)} \left( \frac{\hbar}{\psi_{\pm}}
\frac{\p \psi_{\pm}}{\p x} - p \right).\tag*{\qed}
\end{gather*}
\renewcommand{\qed}{}
\end{proof}

As we will see below, an {\em isomonodromic WKB solution}
such as~\eqref{eq:WKB-IM}
is constructed from just a family of algebraic curves~\eqref{eq:sp-curve}
by {\em the topological recursion}~(\cite{EO07}).
In particular, the f\/irst equation
in \eqref{eq:Lax-scalar} gives a {\em quantization}
of the spectral curve~\eqref{eq:sp-curve}
in the sense of~\cite{DM14, DM14-2}.

\begin{rem}
In the above computation the normalization \eqref{eq:WKB-IM}
is essential. Since $P_{\rm odd}$ is anti-invariant under the
covering involution $\sigma$ as \eqref{eq:parity} and
the integral in~\eqref{eq:WKB-IM} is
def\/ined as a contour integral~\eqref{eq:contour-integral-Podd},
we don't need to take care of the branch point~$v$ in the computation
\begin{gather*}
\int^{x}_{v} \frac{\partial P_{\rm odd}}{\partial t}dx =
\frac{1}{2}\int^{x}_{\sigma(x)}
\frac{\p}{\p x}\left( \frac{P_{\rm odd}}{2(x-q)}\right)dx
= \frac{P_{\rm odd}}{2(x-q)}.
\end{gather*}
\end{rem}

\begin{rem}
We can also construct a WKB-type formal solution
of matrix isomonodromy system \eqref{eq:Lax-matrix}.
Def\/ine
\begin{gather*}
\tilde{\psi}_{\pm}(x,t,\hbar)   =
\frac{ \hbar \frac{d \psi_{\pm}}{dx}(x,t,\hbar) - A_{11}(x,t,\hbar)
\psi_{\pm}(x,t,\hbar)}{A_{12}(x,t,\hbar)}  \\
\hphantom{\tilde{\psi}_{\pm}(x,t,\hbar) }{}
  =
\frac{\hbar P^{(\pm)}(x,t,\hbar) - A_{11}(x,t,\hbar)}
{A_{12}(x,t,\hbar)} \psi_{\pm}(x,t,\hbar).
\end{gather*}
Then, the matrix valued formal series
\begin{gather} \label{eq:WKB-IM-mat}
\Psi(x,t,\hbar)   =
\begin{pmatrix}
\psi_{+}(x,t,\hbar) & \psi_{-}(x,t,\hbar) \\[+.7em]
\tilde{\psi}_{+}(x,t,\hbar) & \tilde{\psi}_{-}(x,t,\hbar)
\end{pmatrix}
\end{gather}
gives a fundamental formal solution of the
isomonodoromy system~\eqref{eq:Lax-matrix}.
\end{rem}

\section{Topological recursion and quantum curve theorem}\label{section3}

In this section we review the Eynard--Orantin's
topological recursion~\cite{EO07} for our
spectral cur\-ve~\eqref{eq:sp-curve},
and formulate our main theorem.

\subsection{Topological recursion}
\label{section:top-recursion}

The topological recursion is an algorithm associating some
dif\/ferential forms $W_{g,n}$ and num\-bers~$F_{g}$
given the following source data:
\begin{itemize}\itemsep=0pt
\item
A plane curve $({\mathcal C}, x, y)$: ${\mathcal C}$ is a
compact Riemann surface, $x, y\colon {\mathcal C} \rightarrow {\mathbb P}^{1}$
are meromorphic functions.
\item
The Bergman kernel $B$:
It is a symmetric dif\/ferential form on ${\mathcal C} \times {\mathcal C}$
with poles of order 2 along the diagonal, and
satisfying some normalization conditions.
\end{itemize}


In our case, ${\mathcal C} = {\mathbb P}^1$ and $x, y$ are
rational functions which parametrize the spectral curve~\eqref{eq:sp-curve}
\begin{gather} \label{eq:par-rep}
x(z) = z^{2} - 2 q_{0}, \qquad
y(z) = 2z(z^{2} - 3q_{0}).
\end{gather}
Here $z$ is a coordinate on ${\mathbb P}^1$.
The Bergman kernel is given by
\begin{gather*}
B(z_1, z_2) = \frac{dz_{1}dz_{2}}{(z_{1}-z_{2})^{2}},
\end{gather*}
since the spectral curve is of genus $0$.
Zeros of~$dx$ are called {\em ramification points} of the
spectral curve~\eqref{eq:par-rep}.
Our spectral curve has only one ramif\/ication point at $z=0$.

The topological recursion for our spectral curve~\eqref{eq:par-rep}
is formulated as follows (see~\cite{EO07} for general case):

\begin{Def}[\protect{\cite[Def\/inition~4.2]{EO07}  (see also~\cite[Section~3]{DM14})}]
The {\em Eynard--Orantin differential}
$W_{g,n}(z_{1},\dots,z_{n})$ of type $(g,n)$
is a meromorphic $n$-dif\/ferential on the $n$-times product
of the spectral curve~\eqref{eq:par-rep} def\/ined by
the following {\em topological recursion relation}:
\begin{itemize}\itemsep=0pt
\item
for $2g-2+n \le 0$:
\begin{gather*}
W_{0,1}(z_{1})   :=   y(z_{1}) dx(z_{1})  = 4z_{1}^{2} \big(z_{1}^{2}-3q_{0}\big)dz_{1}  , \\
W_{0,2}(z_{1},z_{2})   :=   B(z_1, z_2)   = \frac{dz_{1}dz_{2}}{(z_{1}-z_{2})^{2}} ,
\end{gather*}
\item %
for $2g-2+n = 1$:
\begin{gather*}
W_{0,3}(z_{1},z_{2},z_{3})   :=   \frac{1}{2\pi i}
\oint_{\gamma_{0}} K(z,z_{1}) \big[ W_{0,2}(z,z_{2}) W_{0,2}(\bar{z},z_{3})
+ W_{0,2}(z,z_{3}) W_{0,2}(\bar{z},z_{2}) \big], \\
W_{1,1}(z_{1})   :=   \frac{1}{2\pi i}
\oint_{\gamma_{0}} K(z,z_{1}) W_{0,2}(z,\bar{z}),
\end{gather*}
\item %
for $2g-2+n \ge 2$:
\begin{gather}
W_{g,n}(z_{1},\dots,z_{n}) := \frac{1}{2\pi i}
\oint_{\gamma_{0}} K(z,z_{1})   \nonumber \\
 \qquad{}\times  \Biggl[ \sum_{j=2}^{n}
\big( W_{0,2}(z,z_{j}) W_{g,n-1}\big(\bar{z}, z_{[\hat{1},\hat{j}]}\big) +
W_{0,2}(\bar{z},z_{j}) W_{g,n-1}\big(z, z_{[\hat{1},\hat{j}]}\big) \big) \nonumber \\
\qquad{}+ W_{g-1,n+1}\big(z,\bar{z},z_{[\hat{1}]}\big) +
\sum_{\substack{g_{1}+g_{2}=g \\ I \sqcup J = [\hat{1}]}}^{\text{stable}}
W_{g_{1}, |I|+1}(z,z_{I}) W_{g_{2}, |J|+1}(\bar{z}, z_{J}) \Biggr].\label{eq:top-rec}
\end{gather}
\end{itemize}
Here $\gamma_{0}$ is a small cycle (in $z$-plane) which encircles
the ramif\/ication point $z=0$ in the counter-clockwise direction,
$\bar{z} = - z$ is the conjugate of $z$ near the ramif\/ication point,
and the {\em recursion kernel}~$K(z,z_{1})$ is given by
\begin{gather*} %\label{eq:recursion-kernel}
K(z,z_{1}) = - \frac{\omega^{\bar{z}-z}(z_{1})}{2(y(z) - y(\bar{z}))dx(z)}, \qquad
\omega^{\bar{z}-z}(z_{1}) = \int^{\bar{z}}_{z} W_{0,2}( \cdot, z_{1}).
\end{gather*}
Also, we use the index convention
$[\hat{j}] = \{1, \dots, n \}{\setminus}\{j\}$ and so on.
Lastly, the sum in the third line of~\eqref{eq:top-rec} is taken
for indices in the stable range
(i.e., only $W_{g,n}$'s with $2g - 2 + n \ge 1$ appear).
\end{Def}

The explicit form of some of Eynard--Orantin dif\/ferentials
are given as follows
\begin{gather*}
W_{0,3}   =   \frac{1}{12q_{0}z_{1}^{2}z_{2}^{2}z_{3}^{2}}
dz_{1}dz_{2}dz_{3}, \\
W_{0,4}   =   \frac{z_{1}^{2}z_{2}^{2}z_{3}^{2}z_{4}^{2} +
3q_{0} (z_{1}^{2}z_{2}^{2}z_{3}^{2}+z_{2}^{2}z_{3}^{2}z_{4}^{2}+
z_{3}^{2}z_{4}^{2}z_{1}^{2}+
z_{4}^{2}z_{1}^{2}z_{2}^{2})}{144q_{0}^{3}
z_{1}^{4}z_{2}^{4}z_{3}^{4}z_{4}^{4}} dz_{1}dz_{2}dz_{3}dz_{4}, \\
W_{1,1}   =   \frac{z_{1}^{2} + 3q_{0}}
{288 q_{0}^{2} z_{1}^{4}} dz_{1},
\\
W_{1,2}   =   \frac{ 2z_{1}^{4}z_{2}^{4} +
6q_{0}(z_{1}^{4}z_{2}^{2} + z_{1}^{2} z_{2}^{4}) +
3q_{0}^{2} (5z_{1}^{4} + 3 z_{1}^{2}z_{2}^{2} + 5z_{2}^{4})  }
{3456 q_{0}^{4} z_{1}^{6}z_{2}^{6}} dz_{1}dz_{2}, \\
W_{2,1}   =   \frac{ 28 z_{1}^{8} + 84q_{0} z_{1}^{6} + 252q_{0}^{2}z_{1}^{4}
+ 609q_{0}^{3}z_{1}^{2}  + 945 q_{0}^{4}}
{1990656q_{0}^{7} z_{1}^{10}}dz_{1}.
\end{gather*}

Eynard--Orantin dif\/ferentials have
the following properties (see~\cite{EO07}):
\begin{itemize}\itemsep=0pt
\item
As a dif\/ferential form on each variable $z_{i}$,
$W_{g,n}$, for $2g-2+n \ge 1$, is {\em holomorphic}
except for the ramif\/ication point $0$ and may
have a pole at $0$.
\item
$W_{g,n}$ is {\em symmetric}; that is, they are invariant
under any permutation of variables.
\item
For $2g-2+n \ge 1$, $W_{g,n}$
is {\em anti-invariant} under the involution
$z_{i} \mapsto \bar{z}_{i}$ for each variable:
\begin{gather*} %\label{eq:odd-property-of-Wgn}
W_{g,n}(z_{1},\dots,\bar{z}_{j},\dots,z_{n}) = -
W_{g,n}(z_{1},\dots,z_{j},\dots,z_{n}) \qquad
\text{for} \quad j = 1,\dots,n.
\end{gather*}

\item
$W_{g,n}$ is also {\em holomorphic} in $t$ except for
$t = 0$ (i.e., $q_{0} = 0$).
There is a formula for the derivative of $W_{g,n}$
with respect to $t$; see Section~\ref{section:variation-formula}.
\end{itemize}

\subsection{Quantum curve theorem}
\label{section:quantum-curve-conjecture}

In this section we describe our main result which claims
that the scalar isomonodromy system~\eqref{eq:Lax-scalar}
gives a {\em quantum curve}.

\begin{Def}
For $g \ge 0, n \ge 1$ satisfying $2g-2+n \ge 1$,
def\/ine {\em open free energy} of type ($g,n$) by
\begin{gather} \label{eq:open-Fgn}
F_{g,n}(z_{1},\dots,z_{n}) :=
\frac{1}{2^{n}} \int^{z_{1}}_{\bar{z}_{1}}\cdots
\int^{z_{n}}_{\bar{z}_{n}} W_{g,n}(z_{1},\dots,z_{n}).
\end{gather}
\end{Def}

It follows from the def\/inition that open free energies satisfy
\begin{gather*}
d_{z_{1}}\cdots d_{z_{n}} F_{g,n}(z_{1},\dots,z_{n})   =
W_{g,n}(z_{1},\dots,z_{n}), \\
F_{g,n}(z_{1},\dots,\bar{z}_{j}, \dots,z_{n})   =
- F_{g,n}(z_{1},\dots, z_{j}, \dots,z_{n})
\qquad \text{for} \quad j=1,\dots,n.
\end{gather*}
Explicit computation shows that
\begin{gather*}
F_{0,3}(z_{1},z_{2},z_{3})   =   -\frac{1}{12q_{0}z_{1}z_{2}z_{3}}, \\
F_{0,4}(z_{1},z_{2},z_{3},z_{4})   =
\frac{z_{1}^{2}z_{2}^{2}z_{3}^{2}z_{4}^{2} + q_{0}\bigl(
z_{1}^{2}z_{2}^{2}z_{3}^{2} + z_{2}^{2}z_{3}^{2}z_{4}^{2} +
z_{3}^{2}z_{4}^{2}z_{1}^{2} + z_{4}^{2}z_{1}^{2}z_{2}^{2} \bigr)}
{144q_{0}^{3}z_{1}^{3}z_{2}^{3}z_{3}^{3}z_{4}^{3}},
\\
F_{1,1}(z_{1})  =  -\frac{z_{1}^{2}+q_{0}}{288q_{0}^{2} z_{1}^{3}}, \\
%
F_{1,2}(z_{1},z_{2}) =
\frac{2z_{1}^{4}z_{2}^{4} +
2q_{0}\bigl( z_{1}^{4}z_{2}^{2} + z_{1}^{2}z_{2}^{4} \bigr) +
q_{0}^{2}\bigl( 3z_{1}^{4} + z_{1}^{2}z_{2}^{2} + 3z_{2}^{4} \bigr)}
{3456q_{0}^{4}z_{1}^{5}z_{2}^{5}}, \\
%
F_{2,1}(z_{1})  =
-\frac{140z_{1}^{8} + 140 q_{0} z_{1}^6 + 252 q_{0}^{2} z_{1}^4
+ 435 q_{0}^{3} z_{1}^2 + 525 q_{0}^{4}}
{9953280 q_{0}^{7} z_{1}^{9}}.
\end{gather*}
We also introduce functions $\{ S_{m}(x,t) \}_{m \ge 0}$ by
\begin{gather*} %\label{eq:S0-and-S1}
S_{0}(x,t) := \int^{x}_{v} y(z(x'))dx', \qquad
S_{1}(x,t) := -\frac{1}{2} \log \left(\frac{y(z(x))}{2(x-q_{0})}
\right),
\end{gather*}
and for $m \ge 2$
\begin{gather*} %\label{eq:higher-Sm}
S_{m}(x,t) := \sum_{\substack{2g-2+n=m-1 \\ g\ge0, n\ge 1}}
\frac{F_{g,n}(z,\dots,z)}{n!} \biggr|_{z=z(x)},
\end{gather*}
where $z(x) = \sqrt{x+2q_{0}}$ is the inverse function of $x(z)$.
After computations we have
\begin{gather*}
S_{0}(x,t)  =  \frac{4}{5}(x-3q_{0})(x+2q_{0})^{3/2}, \qquad
S_{1}(x,t)  =  - \frac{1}{4}\log(x+2q_{0}), \\
S_{2}(x,t)  =  -\frac{x+7q_{0}}{288q_{0}^{2}(x+2q_{0})^{3/2}}, \qquad
S_{3}(x,t)  =  \frac{2x^{2}+14q_{0}x+35q_{0}^{2}}
{6912q_{0}^{4}(x+2q_{0})^{3}}, \\
S_{4}(x,t)  =  -\frac{140x^{4}+1580q_{0}x^{3}
+ 7476 q_{0}^{2}x^{2}+18739q_{0}^{3}x+23499q_{0}^{4}}
{9953280 q_{0}^{7}(x+2q_{0})^{9/2}}.
\end{gather*}

Our main result is the following.

\begin{thm} \label{conj:2}
